\section{From cross mass to triangle mass}
\label{sec:triangle}

The triangle argument will replace vertices by copies carrying their full
original weights. Total vertex weight can therefore increase. The rule
that makes this operation useful is that each original edge is retained
at most once. We show that this rule preserves the cross-mass bound and
controls edge and triangle masses, so the subsequent argument can proceed
without maximizing weights anew.

Retain the notation $M_G(w)$, $T_G(w)$, $e_w(A,B)$ and $g_x$
from the preceding sections. We write $e_{G,w}$ when the graph must
be specified. We say that $(G,w)$ satisfies
the cross-mass bound with constant $C$ if, for every $x\ge1$ and every
$A,B\subseteq V(G)$,
\begin{equation}
 \label{eq:cross-mass}
 w(A),w(B)\le e^x
 \quad\Longrightarrow\quad
 e_{G,w}(A,B)\le Cg_x(w(A)).
\end{equation}

\begin{theorem}[Triangle mass from cross mass]
\label{thm:weighted-triangles}
For every $C\ge1$ there is a constant $C_{\triangle}(C)$ with the
following property.  If $G$ is a finite simple graph with positive
vertex weights $w$ satisfying \eqref{eq:cross-mass}, then
\[
 T_G(w)\le C_{\triangle}(C)M_G(w).
\]
\end{theorem}

We will repeatedly split vertices so as to reduce neighborhood weights,
while retaining a fixed proportion of triangle weight.  The cross-mass
bound is the information that passes from one step to the next.  We first
verify this preservation and record the effect of deleting edges whose
common-neighbor weight is small.  The proof of
Theorem~\ref{thm:weighted-triangles} concludes after the splitting
construction and its iteration.

\subsection{Preservation and triangle counting}

\begin{lemma}[Preservation under splitting]
\label{lem:invariant-preservation}
Let $(G,w)$ satisfy \eqref{eq:cross-mass} with constant $C$.
Let $H$ be a finite simple graph and let $\pi:V(H)\to V(G)$ satisfy
the following conditions:
\begin{enumerate}
 \item every edge $\{a,b\}$ of $H$ maps to the edge
       $\{\pi(a),\pi(b)\}$ of $G$;
 \item the map $(a,b)\mapsto(\pi(a),\pi(b))$ is injective on
       directed edges of $H$.
\end{enumerate}
Give each $a\in V(H)$ the weight $\widetilde w_a=w_{\pi(a)}$.
Then $(H,\widetilde w)$ satisfies
\eqref{eq:cross-mass} with the same constant $C$.  Moreover,
\begin{gather*}
 M_H(\widetilde w)\le M_G(w),\qquad
 T_H(\widetilde w)\le T_G(w),\\
 \widetilde w(N_H(a))\le w(N_G(\pi(a)))
 \quad\text{for every }a\in V(H).
\end{gather*}
These conclusions are preserved under any finite composition of such maps.
\end{lemma}

\begin{proof}
For $A,B\subseteq V(H)$, write $A_0=\pi(A)$ and $B_0=\pi(B)$.
Duplicating vertices can increase their total weight, so
\[
 w(A_0)\le\widetilde w(A),\qquad
 w(B_0)\le\widetilde w(B).
\]
The directed edges counted by $e_{H,\widetilde w}(A,B)$ map injectively
to directed edges counted by $e_{G,w}(A_0,B_0)$, with their weights
unchanged.  Consequently, if $\widetilde w(A),\widetilde w(B)\le e^x$,
then
\[
 e_{H,\widetilde w}(A,B)
 \le e_{G,w}(A_0,B_0)
 \le Cg_x(w(A_0))
 \le Cg_x(\widetilde w(A)).
\]
This also covers overlapping sets and the case of an empty set.

The directed-edge condition implies injectivity on unordered edges:
two edges with the same unordered image could be oriented to have the
same directed image.  It therefore gives the assertion about $M$.
For a fixed $a$, two different neighbors $b,b'$ cannot have the same
image, since the directed edges $(a,b)$ and $(a,b')$ would then have the
same image.  Thus $\pi$ maps $N_H(a)$ injectively into $N_G(\pi(a))$,
giving the neighborhood bound.

Every triangle of $H$ maps to a triangle of $G$, since a homomorphism
to a simple graph is injective on the vertices of a clique.  Distinct
triangles cannot have the same image: each of the three image edges
has at most one preimage edge, and those three preimage edges determine
the triangle.  Triangle weights are unchanged, proving the assertion
about $T$.  Finally, compositions preserve vertex weights,
adjacency and injectivity on directed edges.
\end{proof}

Edge deletion and vertex deletion are special cases of
Lemma~\ref{lem:invariant-preservation}, using the inclusion map.
The splitting operation used below is another special case: replace a
vertex $u$ by copies indexed by groups of its incident edges, give every
copy weight $w_u$, and retain each original edge at most once, between
its assigned endpoint copies.  No assumption about optimization is
required after these operations.

\begin{lemma}[Triangle normalization and peeling]
\label{lem:triangle-peeling}
For a positive-weighted finite simple graph $G$, put
\[
 s_u=w(N_G(u)),\qquad
 c_{uv}=w(N_G(u)\cap N_G(v)),\qquad
 h_u=\sum_{v\in N_G(u)}w_vc_{uv}.
\]
Thus $h_u=2M_{G[N_G(u)]}(w)$ is twice the edge mass inside the
neighborhood of $u$.
Then
\begin{equation}
 \label{eq:triangle-normalizations}
 \sum_{\{u,v\}\in E(G)}w_uw_vc_{uv}=3T_G(w),
 \qquad
 \sum_{u\in V(G)}w_uh_u=6T_G(w).
\end{equation}
In particular, if every neighborhood has weight at most $L$, then
\begin{equation}
 \label{eq:terminal-triangle-bound}
 T_G(w)\le \frac{L}{3}M_G(w).
\end{equation}

For every $\tau>0$, there is a spanning subgraph $G'$ obtained by
successive edge deletions such that every edge of $G'$ has common-neighbor
weight at least $\tau$ and
\begin{equation}
 \label{eq:peeling-loss}
 T_{G'}(w)\ge T_G(w)-\tau M_G(w).
\end{equation}
For each nonisolated vertex of $G'$, its neighborhood and common-neighbor
weights, computed in $G'$, satisfy
\[
 s_u\ge\tau,\qquad h_u\ge\tau s_u.
\]
If $G$ satisfies \eqref{eq:cross-mass}, then so does $G'$,
with the same constant.
\end{lemma}

\begin{proof}
Each triangle contributes its weight once for each of its three edges
to the first sum in \eqref{eq:triangle-normalizations}, and once for
each of its six oriented edges to the second.  Since $c_{uv}\le L$,
the first identity also proves \eqref{eq:terminal-triangle-bound}.

Starting with $G$, delete any edge whose common-neighbor weight in the
current graph is less than $\tau$, and continue until no such edge
remains.  The procedure terminates because $G$ is finite.  If
$\{u_i,v_i\}$ is the $i$th deleted edge and $c_i$ its common-neighbor
weight immediately before deletion, the triangle weight lost at that
step is exactly $w_{u_i}w_{v_i}c_i$.  Telescoping gives
\[
 T_G(w)-T_{G'}(w)
 =\sum_iw_{u_i}w_{v_i}c_i
 \le\tau\sum_iw_{u_i}w_{v_i}
 \le\tau M_G(w).
\]
This counts each destroyed triangle at its first deleted edge.
The defining stopping rule gives $c_{uv}\ge\tau$ on every remaining
edge.  If $u$ is nonisolated, choose a remaining neighbor $v$ to obtain
$s_u\ge c_{uv}\ge\tau$; summing $w_vc_{uv}\ge\tau w_v$ over the
remaining neighbors gives $h_u\ge\tau s_u$.  The last assertion follows
from Lemma~\ref{lem:invariant-preservation}.
\end{proof}
